% =============================================================================
%  Capítulo 3 — Diagramas de Referência (Parte III)
%  Diagramas do projeto; os PDFs vetoriais ficam em img/diagramas/ (cores do guia).
% =============================================================================
\chapter{Diagramas de Referência}\label{cap:ref}

Esta parte reúne os diagramas técnicos do sistema, para quem precisa ir além da visão
geral da Parte~I: o que cada perfil pode fazer (casos de uso), como o agente de IA e a consulta funcionam por dentro, como os componentes se
organizam, onde cada parte roda, como o código se divide em pacotes, o que acontece em cada fluxo importante e como os dados se relacionam. Os desenhos são vetoriais; em uma tela, é
possível ampliar qualquer trecho.

% Gerado por gen.py (casos de uso). Não editar à mão: regenerar a partir de ucs.py.
\newtcolorbox{ucard}[2]{enhanced,breakable,colback=guiapapel,colframe=guiaborda,boxrule=.6pt,arc=0pt,
  left=8pt,right=8pt,top=5pt,bottom=2pt,colbacktitle=guiavermelho,coltitle=white,fonttitle=\bfseries\small,
  toptitle=3pt,bottomtitle=3pt,title={Caso de uso: #1\enskip---\enskip #2},before skip=10pt,after skip=12pt}
\newcommand{\ucrow}[2]{{\color{guiavermelho}\bfseries\small #1:}\par\vspace{1pt}{\small #2\par}\vspace{3pt}{\color{guiaborda}\hrule height .4pt}\vspace{4pt}}
\section{Casos de uso}\label{sec:casos-de-uso}

O diagrama da \autoref{fig:casos-de-uso} segue a notação da UML. O retângulo é a fronteira do sistema, o
Hanaro; as elipses são os casos de uso; os bonecos à esquerda são os atores humanos e os da direita são
sistemas externos que participam de um caso. Uma linha contínua liga o ator ao caso de que participa. A seta
com triângulo vazado é generalização: o \emph{Gestor} faz tudo o que o \emph{Analista} faz, e o
\emph{Superusuário} tudo o que o \emph{Administrador} faz; \emph{Analista} e \emph{Administrador} são
tipos de \emph{Usuário autenticado}, ator abstrato (em itálico) que reúne o que todas as contas fazem. A seta
tracejada «include» indica que um caso sempre usa o outro; «extend», que um caso acrescenta um
comportamento opcional ao caso que estende.

Os casos de uso de cada ator foram levantados das telas, das rotas e das regras de autorização do sistema.
Cada caso tem, em seguida, um card de especificação com referências, descrição, atores, pré-condições,
garantia de sucesso, requisitos especiais e os fluxos básico e alternativo.

\begin{figure}[H]
\centering
\includegraphics[width=\textwidth,height=\dimexpr\textheight-4.2cm\relax,keepaspectratio]{img/diagramas/casos-de-uso}
\caption{Diagrama de casos de uso do Hanaro.}\label{fig:casos-de-uso}
\end{figure}

\Needspace{14\baselineskip}
\subsection{Acesso, execuções e entrada de dados}

\Needspace{9\baselineskip}
\begin{ucard}{UC01}{Autenticar-se}
\ucrow{Referências}{Tela: Tela de acesso (/login). Fluxo de sequência: \autoref{fig:seq01}.}
\ucrow{Descrição geral}{O caso de uso inicia-se quando uma pessoa com conta no Hanaro informa usuário e senha para entrar no sistema.}
\ucrow{Atores}{Usuário autenticado (Analista, Gestor ou Administrador).}
\ucrow{Pré-condições}{Conta provisionada pelo administrador.}
\ucrow{Garantia de sucesso (pós-condições)}{Sessão aberta. O navegador guarda só o identificador da sessão, em cookie \texttt{HttpOnly}, e recebe o token CSRF usado nas gravações.}
\ucrow{Requisitos especiais}{O único acesso é por usuário e senha; não há login social. Há limite de requisições por IP.}
\ucrow{Fluxo básico}{\begin{enumerate}[nosep,leftmargin=*,label=\arabic*.]
\item A pessoa abre a tela de acesso e informa usuário e senha.
\item O sistema valida as credenciais.
\item O sistema cria a sessão e devolve o token CSRF.
\item O sistema leva a pessoa à tela inicial do seu perfil (Execuções para o administrador; Dashboard para os demais).
\end{enumerate}}
\ucrow{Fluxo alternativo}{\begin{enumerate}[nosep,leftmargin=*,label=A\arabic*.]
\item Credenciais inválidas: o sistema informa o erro e permanece na tela de acesso. Retorna ao passo 1.
\item Sessão expirada durante o uso: o sistema abre a janela de acesso sobre a tela atual; ao entrar de novo, a pessoa continua de onde parou.
\item Limite de requisições excedido: o sistema recusa o pedido e orienta tentar mais tarde.
\end{enumerate}}
\end{ucard}

\Needspace{9\baselineskip}
\begin{ucard}{UC02}{Gerenciar perfil e sessão}
\ucrow{Referências}{Tela: Perfil (/perfil). Fluxo de sequência: \autoref{fig:seq01}.}
\ucrow{Descrição geral}{O caso de uso inicia-se quando a pessoa autenticada quer consultar ou alterar seus dados de contato e suas preferências, ou encerrar a sessão.}
\ucrow{Atores}{Usuário autenticado.}
\ucrow{Pré-condições}{Sessão ativa.}
\ucrow{Garantia de sucesso (pós-condições)}{Telefone e preferências atualizados e persistidos, ou sessão encerrada.}
\ucrow{Requisitos especiais}{A gravação exige o token CSRF da sessão.}
\ucrow{Fluxo básico}{\begin{enumerate}[nosep,leftmargin=*,label=\arabic*.]
\item A pessoa abre a tela de Perfil.
\item O sistema mostra o telefone e as preferências atuais.
\item A pessoa altera os campos desejados e confirma.
\item O sistema valida e grava as alterações.
\item O sistema confirma a gravação.
\end{enumerate}}
\ucrow{Fluxo alternativo}{\begin{enumerate}[nosep,leftmargin=*,label=A\arabic*.]
\item Valor inválido (por exemplo, telefone fora do formato): o sistema indica o campo e não grava. Retorna ao passo 3.
\item Encerrar sessão: a pessoa escolhe sair; o sistema invalida a sessão e volta à tela de acesso.
\end{enumerate}}
\end{ucard}

\Needspace{9\baselineskip}
\begin{ucard}{UC03}{Solicitar suporte}
\ucrow{Referências}{Tela: Janela Fale com o suporte.}
\ucrow{Descrição geral}{O caso de uso inicia-se quando a pessoa autenticada precisa relatar uma dúvida, um problema ou uma sugestão à equipe de suporte.}
\ucrow{Atores}{Usuário autenticado.}
\ucrow{Pré-condições}{Sessão ativa.}
\ucrow{Garantia de sucesso (pós-condições)}{Solicitação registrada, com assunto, mensagem e dados de quem enviou, disponível para o administrador (UC23).}
\ucrow{Requisitos especiais}{Nenhum.}
\ucrow{Fluxo básico}{\begin{enumerate}[nosep,leftmargin=*,label=\arabic*.]
\item A pessoa abre a janela Fale com o suporte.
\item A pessoa informa o assunto e descreve o pedido.
\item A pessoa envia a solicitação.
\item O sistema valida os campos e registra a solicitação.
\item O sistema confirma o envio.
\end{enumerate}}
\ucrow{Fluxo alternativo}{\begin{enumerate}[nosep,leftmargin=*,label=A\arabic*.]
\item Assunto ou mensagem vazios: o sistema indica o erro e não envia. Retorna ao passo 2.
\item Falha no envio: o sistema avisa e permite tentar de novo; o texto digitado é mantido.
\end{enumerate}}
\end{ucard}

\Needspace{9\baselineskip}
\begin{ucard}{UC04}{Acompanhar execuções}
\ucrow{Referências}{Tela: Execuções (/execucoes). Fluxo de sequência: \autoref{fig:seq03}.}
\ucrow{Descrição geral}{O caso de uso inicia-se quando a pessoa quer saber se as cargas de scrap chegaram ao Hanaro e como terminaram.}
\ucrow{Atores}{Usuário autenticado.}
\ucrow{Pré-condições}{Sessão ativa.}
\ucrow{Garantia de sucesso (pós-condições)}{A pessoa conhece a situação de cada execução, suas contagens e seus erros.}
\ucrow{Requisitos especiais}{A lista é paginada e pode ser buscada.}
\ucrow{Fluxo básico}{\begin{enumerate}[nosep,leftmargin=*,label=\arabic*.]
\item A pessoa abre a tela de Execuções.
\item O sistema lista as execuções recentes, com situação e horário.
\item A pessoa busca ou ordena a lista, se quiser.
\item A pessoa abre o detalhe de uma execução.
\item O sistema mostra o resultado, as contagens e os erros dessa execução.
\end{enumerate}}
\ucrow{Fluxo alternativo}{\begin{enumerate}[nosep,leftmargin=*,label=A\arabic*.]
\item Nenhuma execução registrada: o sistema mostra o estado vazio.
\item Falha ao carregar: o sistema mostra o erro e oferece tentar novamente.
\end{enumerate}}
\end{ucard}

\Needspace{9\baselineskip}
\begin{ucard}{UC05}{Enviar arquivo de scrap manualmente}
\ucrow{Referências}{Tela: Execuções, janela de envio manual. Fluxo de sequência: \autoref{fig:seq03}.}
\ucrow{Descrição geral}{O caso de uso estende Acompanhar execuções (UC04) e inicia-se quando a pessoa precisa carregar à mão um arquivo exportado do GERP.}
\ucrow{Atores}{Usuário autenticado.}
\ucrow{Pré-condições}{\begin{enumerate}[nosep,leftmargin=*,label=\arabic*.]
\item Arquivo exportado do GERP em mãos.
\item Sessão ativa.
\end{enumerate}}
\ucrow{Garantia de sucesso (pós-condições)}{Execução criada e enfileirada; ao final, os dados do arquivo entram na base de scrap e na projeção do Dashboard.}
\ucrow{Requisitos especiais}{A cotação usada é a do dia de chegada informado (hoje, em Manaus, por padrão). O mesmo arquivo reenviado na mesma janela não duplica registros.}
\ucrow{Fluxo básico}{\begin{enumerate}[nosep,leftmargin=*,label=\arabic*.]
\item A pessoa escolhe enviar um arquivo na tela de Execuções.
\item A pessoa seleciona o arquivo e confirma o dia de chegada.
\item O sistema valida e converte o arquivo na carga canônica.
\item O sistema cria a execução e a enfileira.
\item O worker reconcilia e grava os dados.
\item A pessoa acompanha a execução até ser concluída (UC04).
\end{enumerate}}
\ucrow{Fluxo alternativo}{\begin{enumerate}[nosep,leftmargin=*,label=A\arabic*.]
\item Arquivo inválido: o sistema lista os erros, nada é gravado e a pessoa pode corrigir e reenviar. Retorna ao passo 2.
\item Sem cotação no dia informado: vale a última cotação disponível.
\item Lote com falha: a execução é marcada com falha e o último estado válido é mantido.
\end{enumerate}}
\end{ucard}

\Needspace{9\baselineskip}
\begin{ucard}{UC06}{Ingerir carga de scrap do GERP}
\ucrow{Referências}{Tela: API de ingestão (sem tela). Fluxo de sequência: \autoref{fig:seq02}.}
\ucrow{Descrição geral}{O caso de uso inicia-se quando o robô do GERP envia ao Hanaro a carga de scrap convertida no formato canônico.}
\ucrow{Atores}{Robô do GERP.}
\ucrow{Pré-condições}{Chave de API válida, enviada no cabeçalho \texttt{X-API-Key}.}
\ucrow{Garantia de sucesso (pós-condições)}{Carga enfileirada e processada: transações, projeção do Dashboard e snapshot ativo gravados no mesmo commit; execução marcada como concluída.}
\ucrow{Requisitos especiais}{A chave só permite criar registros de scrap. A API valida a forma e responde \texttt{202}; a reconciliação é do worker, por partição (organização e dia).}
\ucrow{Fluxo básico}{\begin{enumerate}[nosep,leftmargin=*,label=\arabic*.]
\item O robô envia a carga à API com a chave de API.
\item A API confere a chave e a permissão de criação.
\item A API valida a forma da carga.
\item A API enfileira a tarefa e responde 202 com o identificador.
\item O worker valida contagens, totais, período e USD.
\item O worker reconcilia as partições cobertas e grava transações, projeção e snapshot.
\item O worker liga o resultado à execução e a marca como concluída.
\end{enumerate}}
\ucrow{Fluxo alternativo}{\begin{enumerate}[nosep,leftmargin=*,label=A\arabic*.]
\item Chave inválida ou ausente: a API recusa o pedido (401/403).
\item Carga fora do formato: a API recusa e informa o erro.
\item Lote com falha: a execução é marcada com falha e o último estado válido é mantido.
\item Mesmo arquivo na mesma janela: o processamento é idempotente e não duplica.
\end{enumerate}}
\end{ucard}

\Needspace{14\baselineskip}
\subsection{Análise, assistente e relatórios}

\Needspace{9\baselineskip}
\begin{ucard}{UC07}{Consultar dashboard}
\ucrow{Referências}{Tela: Dashboard (/dashboard) e modo TV. Fluxo de sequência: \autoref{fig:seq05}.}
\ucrow{Descrição geral}{O caso de uso inicia-se quando a pessoa quer ver os indicadores de scrap consolidados.}
\ucrow{Atores}{Analista (o Gestor herda o caso de uso).}
\ucrow{Pré-condições}{Existência de dados ingeridos (UC06 ou UC05).}
\ucrow{Garantia de sucesso (pós-condições)}{Indicadores exibidos conforme os filtros escolhidos.}
\ucrow{Requisitos especiais}{O Dashboard lê uma projeção já agregada, não as transações; a chave do cache inclui a revisão da projeção. O modo TV é público, somente leitura e sem dados pessoais.}
\ucrow{Fluxo básico}{\begin{enumerate}[nosep,leftmargin=*,label=\arabic*.]
\item A pessoa abre o Dashboard.
\item O sistema busca a projeção agregada, usando o cache quando houver.
\item O sistema exibe os indicadores.
\item A pessoa ajusta filtros de período e de escopo.
\item O sistema atualiza os indicadores.
\end{enumerate}}
\ucrow{Fluxo alternativo}{\begin{enumerate}[nosep,leftmargin=*,label=A\arabic*.]
\item Cache indisponível: o sistema lê direto do banco e a resposta apenas demora mais.
\item Sem dados no período: o sistema mostra o estado vazio.
\item Nova ingestão ou mudança de meta: a revisão da projeção muda e o cache antigo deixa de ser usado.
\end{enumerate}}
\end{ucard}

\Needspace{9\baselineskip}
\begin{ucard}{UC08}{Consultar base de scrap}
\ucrow{Referências}{Tela: Base de scrap (/base-de-scrap). Fluxo de sequência: \autoref{fig:seq06}.}
\ucrow{Descrição geral}{O caso de uso inicia-se quando a pessoa quer localizar ocorrências de scrap para conferir, analisar ou usar em relatórios.}
\ucrow{Atores}{Analista.}
\ucrow{Pré-condições}{Sessão ativa.}
\ucrow{Garantia de sucesso (pós-condições)}{Lista de ocorrências exibida conforme os filtros; a pessoa pode abrir uma ocorrência.}
\ucrow{Requisitos especiais}{O acesso exige sessão. Os filtros podem ser salvos como modelos.}
\ucrow{Fluxo básico}{\begin{enumerate}[nosep,leftmargin=*,label=\arabic*.]
\item A pessoa abre a Base de scrap.
\item O sistema lista as ocorrências.
\item A pessoa aplica filtros, busca ou um modelo salvo.
\item O sistema atualiza a lista.
\item A pessoa abre uma ocorrência para ver os detalhes.
\end{enumerate}}
\ucrow{Fluxo alternativo}{\begin{enumerate}[nosep,leftmargin=*,label=A\arabic*.]
\item Nenhum resultado: o sistema informa e sugere limpar os filtros.
\item Salvar modelo: a pessoa nomeia o conjunto de filtros atual e o sistema o guarda para reutilização.
\end{enumerate}}
\end{ucard}

\Needspace{9\baselineskip}
\begin{ucard}{UC09}{Analisar ocorrência de scrap}
\ucrow{Referências}{Tela: Base de scrap, tela de revisão. Fluxo de sequência: \autoref{fig:seq06}.}
\ucrow{Descrição geral}{O caso de uso inicia-se quando a analista trata uma ocorrência de scrap e registra a análise do caso.}
\ucrow{Atores}{Analista.}
\ucrow{Pré-condições}{Ocorrência existente na base de scrap (UC08).}
\ucrow{Garantia de sucesso (pós-condições)}{Nova versão da análise do caso gravada; ao ser marcada como revisada, fica disponível para relatórios.}
\ucrow{Requisitos especiais}{Cada gravação cria uma versão; versões anteriores são preservadas. As análises aprovadas servem de exemplo à memória do assistente.}
\ucrow{Fluxo básico}{\begin{enumerate}[nosep,leftmargin=*,label=\arabic*.]
\item A pessoa localiza a ocorrência na Base de scrap (UC08, incluído).
\item A pessoa abre a tela de revisão.
\item O sistema mostra os fatos da ocorrência e a análise atual, se houver.
\item A pessoa preenche ou altera os campos da análise.
\item A pessoa salva.
\item O sistema grava uma nova versão da análise do caso.
\item A pessoa marca a análise como revisada.
\end{enumerate}}
\ucrow{Fluxo alternativo}{\begin{enumerate}[nosep,leftmargin=*,label=A\arabic*.]
\item Campo obrigatório vazio: o sistema indica e não grava. Retorna ao passo 4.
\item Existe pré-análise do assistente: a pessoa a revisa (UC10, extensão) e pode aproveitar as propostas.
\end{enumerate}}
\end{ucard}

\Needspace{9\baselineskip}
\begin{ucard}{UC10}{Revisar pré-análise do assistente}
\ucrow{Referências}{Tela: Assistente (/assistente). Fluxo de sequência: \autoref{fig:seq08}.}
\ucrow{Descrição geral}{O caso de uso estende Analisar ocorrência (UC09) e inicia-se quando o assistente deixa uma pré-análise pronta para uma ocorrência.}
\ucrow{Atores}{Analista; Servidor de IA (ator secundário).}
\ucrow{Pré-condições}{\begin{enumerate}[nosep,leftmargin=*,label=\arabic*.]
\item Regra de alerta com pré-análise ligada disparou para a ocorrência.
\item Agente de IA disponível no momento do preparo.
\end{enumerate}}
\ucrow{Garantia de sucesso (pós-condições)}{Decisão da pessoa (aceitar, ajustar ou rejeitar) registrada; campos e tarefas aceitos aplicados à análise.}
\ucrow{Requisitos especiais}{A IA propõe e a pessoa decide. As propostas levam selo de origem e vêm com casos parecidos e fontes. A rejeição pede um motivo.}
\ucrow{Fluxo básico}{\begin{enumerate}[nosep,leftmargin=*,label=\arabic*.]
\item A pessoa abre a lista do Assistente e escolhe uma pré-análise.
\item O sistema mostra os fatos, os casos parecidos e as propostas de campos e de tarefas.
\item A pessoa lê as propostas e confere as fontes.
\item A pessoa aceita, ajusta ou rejeita cada proposta.
\item O sistema aplica o que foi aceito à análise e registra a decisão.
\end{enumerate}}
\ucrow{Fluxo alternativo}{\begin{enumerate}[nosep,leftmargin=*,label=A\arabic*.]
\item Agente indisponível: não há pré-análise; só ela atrasa e a pessoa segue com a análise manual (UC09).
\item Rejeição: a pessoa informa o motivo e o sistema registra a decisão sem alterar a análise.
\item Pré-análise sem casos parecidos: o sistema informa e mostra apenas as propostas.
\end{enumerate}}
\end{ucard}

\Needspace{9\baselineskip}
\begin{ucard}{UC11}{Consultar o assistente}
\ucrow{Referências}{Tela: Assistente, aba Consulta. Fluxo de sequência: \autoref{fig:seq09}.}
\ucrow{Descrição geral}{O caso de uso inicia-se quando a pessoa faz ao assistente uma pergunta em texto livre sobre casos já tratados.}
\ucrow{Atores}{Analista; Servidor de IA (ator secundário).}
\ucrow{Pré-condições}{\begin{enumerate}[nosep,leftmargin=*,label=\arabic*.]
\item Recurso de Consulta ligado.
\item Sessão ativa.
\end{enumerate}}
\ucrow{Garantia de sucesso (pós-condições)}{Resposta exibida com as fontes; conversa guardada no histórico da pessoa.}
\ucrow{Requisitos especiais}{A resposta usa só análises já aprovadas e sempre cita as fontes.}
\ucrow{Fluxo básico}{\begin{enumerate}[nosep,leftmargin=*,label=\arabic*.]
\item A pessoa abre a aba Consulta e escreve a pergunta.
\item O sistema envia a pergunta ao servidor de IA.
\item O servidor busca nas análises aprovadas e elabora a resposta.
\item O sistema exibe a resposta com as fontes.
\item O sistema grava a conversa no histórico.
\end{enumerate}}
\ucrow{Fluxo alternativo}{\begin{enumerate}[nosep,leftmargin=*,label=A\arabic*.]
\item Sem base suficiente: o assistente informa que não encontrou casos que sustentem a resposta.
\item Servidor de IA indisponível: o sistema mostra a falha e permite repetir a pergunta.
\item Reabrir conversa: a pessoa escolhe uma conversa do histórico e continua dela.
\end{enumerate}}
\end{ucard}

\Needspace{9\baselineskip}
\begin{ucard}{UC12}{Elaborar relatório}
\ucrow{Referências}{Tela: Relatórios (/relatorios). Fluxo de sequência: \autoref{fig:seq10}.}
\ucrow{Descrição geral}{O caso de uso inicia-se quando a pessoa monta um relatório sobre um evento que gerou scrap relevante.}
\ucrow{Atores}{Analista.}
\ucrow{Pré-condições}{Ao menos uma análise revisada (UC09).}
\ucrow{Garantia de sucesso (pós-condições)}{Relatório salvo como rascunho, com as análises vinculadas e o conteúdo editado.}
\ucrow{Requisitos especiais}{Só análises revisadas podem ser vinculadas. A prévia mostra o documento como será emitido.}
\ucrow{Fluxo básico}{\begin{enumerate}[nosep,leftmargin=*,label=\arabic*.]
\item A pessoa cria um novo relatório.
\item O sistema lista as análises revisadas candidatas.
\item A pessoa vincula as análises e preenche o conteúdo.
\item O sistema exibe a prévia do relatório.
\item A pessoa salva o rascunho.
\end{enumerate}}
\ucrow{Fluxo alternativo}{\begin{enumerate}[nosep,leftmargin=*,label=A\arabic*.]
\item Nenhuma análise revisada: o sistema orienta a revisar uma análise antes (UC09).
\item Excluir relatório: o sistema pede confirmação e remove as ações e vínculos; as análises permanecem na base de scrap.
\end{enumerate}}
\end{ucard}

\Needspace{9\baselineskip}
\begin{ucard}{UC13}{Publicar relatório}
\ucrow{Referências}{Tela: Relatórios, publicação. Fluxo de sequência: \autoref{fig:seq10}.}
\ucrow{Descrição geral}{O caso de uso inicia-se quando a pessoa decide congelar o relatório numa versão oficial.}
\ucrow{Atores}{Analista.}
\ucrow{Pré-condições}{Relatório em rascunho (UC12).}
\ucrow{Garantia de sucesso (pós-condições)}{Versão numerada e imutável criada, selada por um snapshot dos dados do momento.}
\ucrow{Requisitos especiais}{A publicação aceita uma chave de idempotência: um clique repetido não gera duas versões.}
\ucrow{Fluxo básico}{\begin{enumerate}[nosep,leftmargin=*,label=\arabic*.]
\item A pessoa pede a publicação do relatório.
\item O sistema grava o snapshot dos dados com hash.
\item O sistema cria a versão publicada numerada.
\item O sistema confirma a publicação.
\end{enumerate}}
\ucrow{Fluxo alternativo}{\begin{enumerate}[nosep,leftmargin=*,label=A\arabic*.]
\item Pedido repetido com a mesma chave: o sistema devolve a versão já criada.
\item Falha ao gravar: nenhuma versão é criada e a pessoa pode tentar de novo.
\end{enumerate}}
\end{ucard}

\Needspace{9\baselineskip}
\begin{ucard}{UC14}{Exportar relatório}
\ucrow{Referências}{Tela: Relatórios, exportação. Fluxo de sequência: \autoref{fig:seq11}.}
\ucrow{Descrição geral}{O caso de uso inicia-se quando a pessoa precisa do relatório publicado como arquivo para circular.}
\ucrow{Atores}{Analista.}
\ucrow{Pré-condições}{Relatório com ao menos uma versão publicada (UC13).}
\ucrow{Garantia de sucesso (pós-condições)}{Arquivo gerado e disponível para download.}
\ucrow{Requisitos especiais}{A exportação é assíncrona: o pedido vira um job com reserva (lease), para que só um worker o processe e para recuperar o que ficou parado.}
\ucrow{Fluxo básico}{\begin{enumerate}[nosep,leftmargin=*,label=\arabic*.]
\item A pessoa escolhe a versão e pede a exportação.
\item O sistema cria o job e responde de imediato.
\item O worker renderiza o arquivo.
\item A pessoa acompanha o andamento na tela.
\item O sistema informa que o arquivo está pronto e a pessoa o baixa.
\end{enumerate}}
\ucrow{Fluxo alternativo}{\begin{enumerate}[nosep,leftmargin=*,label=A\arabic*.]
\item Falha na geração: o job é marcado com falha, a tela mostra o erro e a pessoa pode pedir de novo.
\item Tradução pedida: a pessoa escolhe emitir em inglês (UC15, extensão).
\end{enumerate}}
\end{ucard}

\Needspace{9\baselineskip}
\begin{ucard}{UC15}{Traduzir relatório para o inglês}
\ucrow{Referências}{Tela: Relatórios, exportação em inglês. Fluxo de sequência: \autoref{fig:seq12}.}
\ucrow{Descrição geral}{O caso de uso estende Exportar relatório (UC14) e inicia-se quando a pessoa pede o documento em inglês.}
\ucrow{Atores}{Analista; Serviço de Tradução (ator secundário).}
\ucrow{Pré-condições}{\begin{enumerate}[nosep,leftmargin=*,label=\arabic*.]
\item Tradução ativada e serviço configurado.
\item Exportação pedida (UC14).
\end{enumerate}}
\ucrow{Garantia de sucesso (pós-condições)}{Texto do relatório traduzido do português para o inglês e incluído no arquivo exportado.}
\ucrow{Requisitos especiais}{Números, códigos e marcas são protegidos antes de traduzir. Só um job fala com o serviço por vez; os demais ficam na fila e a tela mostra a posição.}
\ucrow{Fluxo básico}{\begin{enumerate}[nosep,leftmargin=*,label=\arabic*.]
\item A pessoa marca a opção de emitir em inglês ao exportar.
\item O sistema protege números e códigos e mascara o texto.
\item O sistema envia cada campo ao serviço de tradução.
\item O serviço devolve o texto em inglês.
\item O sistema restaura os trechos protegidos e segue a exportação.
\end{enumerate}}
\ucrow{Fluxo alternativo}{\begin{enumerate}[nosep,leftmargin=*,label=A\arabic*.]
\item Serviço indisponível: o sistema marca a falha da tradução e informa a pessoa; o restante da exportação não é perdido.
\item Fila ocupada: o job aguarda e a tela mostra a posição.
\end{enumerate}}
\end{ucard}

\Needspace{14\baselineskip}
\subsection{Planos, alertas e configurações}

\Needspace{9\baselineskip}
\begin{ucard}{UC16}{Gerenciar plano de ação e tarefas}
\ucrow{Referências}{Tela: Planos de ação (/planos-de-acao).}
\ucrow{Descrição geral}{O caso de uso inicia-se quando a pessoa organiza as ações corretivas e as tarefas que decorrem de um relatório.}
\ucrow{Atores}{Analista.}
\ucrow{Pré-condições}{Relatório existente (UC12).}
\ucrow{Garantia de sucesso (pós-condições)}{Plano, ações e tarefas gravados, com responsáveis e situação.}
\ucrow{Requisitos especiais}{O plano acompanha o relatório e aparece nele ao ser publicado.}
\ucrow{Fluxo básico}{\begin{enumerate}[nosep,leftmargin=*,label=\arabic*.]
\item A pessoa abre os Planos de ação e escolhe um relatório.
\item O sistema mostra o plano e suas tarefas.
\item A pessoa cria ou edita ações e tarefas, com responsável e prazo.
\item O sistema grava as alterações.
\item A pessoa acompanha a situação (aberto, em andamento, fechado).
\end{enumerate}}
\ucrow{Fluxo alternativo}{\begin{enumerate}[nosep,leftmargin=*,label=A\arabic*.]
\item Remover ação: o sistema pede confirmação e a remove do plano.
\item Tarefa com evidência: a pessoa anexa a evidência e o sistema a associa à tarefa.
\end{enumerate}}
\end{ucard}

\Needspace{9\baselineskip}
\begin{ucard}{UC17}{Tratar alertas}
\ucrow{Referências}{Tela: Alertas (/alertas). Fluxo de sequência: \autoref{fig:seq07}.}
\ucrow{Descrição geral}{O caso de uso inicia-se quando a pessoa precisa ver e tratar os alertas gerados pelas regras.}
\ucrow{Atores}{Analista.}
\ucrow{Pré-condições}{Ao menos uma regra de alerta ativa (UC18) cujo evento tenha ocorrido.}
\ucrow{Garantia de sucesso (pós-condições)}{Alertas lidos e tratados; o contador de não lidos é atualizado.}
\ucrow{Requisitos especiais}{O alerta interno é gerado pelo worker; o envio ao Teams, quando configurado, segue o mesmo evento.}
\ucrow{Fluxo básico}{\begin{enumerate}[nosep,leftmargin=*,label=\arabic*.]
\item A pessoa abre a tela de Alertas.
\item O sistema lista os alertas, com o contador de não lidos.
\item A pessoa abre um alerta e lê o evento.
\item A pessoa registra o tratamento.
\item O sistema atualiza a situação do alerta.
\end{enumerate}}
\ucrow{Fluxo alternativo}{\begin{enumerate}[nosep,leftmargin=*,label=A\arabic*.]
\item Sem alertas: o sistema mostra o estado vazio.
\item Alerta com pré-análise: a pessoa segue para a revisão do assistente (UC10).
\end{enumerate}}
\end{ucard}

\Needspace{9\baselineskip}
\begin{ucard}{UC18}{Configurar regras de alerta}
\ucrow{Referências}{Tela: Alertas, regras. Fluxo de sequência: \autoref{fig:seq07}.}
\ucrow{Descrição geral}{O caso de uso inicia-se quando a pessoa define quando o Hanaro deve gerar um alerta.}
\ucrow{Atores}{Analista.}
\ucrow{Pré-condições}{Sessão ativa.}
\ucrow{Garantia de sucesso (pós-condições)}{Regra gravada e avaliada pelo worker a cada ciclo.}
\ucrow{Requisitos especiais}{A regra tem nível e pode ligar a pré-análise do assistente.}
\ucrow{Fluxo básico}{\begin{enumerate}[nosep,leftmargin=*,label=\arabic*.]
\item A pessoa abre as regras na tela de Alertas.
\item A pessoa cria uma regra ou edita uma existente.
\item A pessoa define a condição, o nível e a entrega.
\item O sistema valida e grava a regra.
\item O worker passa a avaliá-la a cada ciclo.
\end{enumerate}}
\ucrow{Fluxo alternativo}{\begin{enumerate}[nosep,leftmargin=*,label=A\arabic*.]
\item Regra inválida: o sistema indica o campo e não grava.
\item Desativar ou excluir regra: o sistema deixa de avaliá-la; alertas já gerados permanecem.
\end{enumerate}}
\end{ucard}

\Needspace{9\baselineskip}
\begin{ucard}{UC19}{Configurar metas e parâmetros}
\ucrow{Referências}{Tela: Configurações (/configuracoes).}
\ucrow{Descrição geral}{O caso de uso inicia-se quando a pessoa ajusta os parâmetros que alimentam os indicadores e os relatórios.}
\ucrow{Atores}{Analista.}
\ucrow{Pré-condições}{Sessão ativa.}
\ucrow{Garantia de sucesso (pós-condições)}{Parâmetros gravados; a mudança de meta troca a revisão da projeção do Dashboard.}
\ucrow{Requisitos especiais}{Cada janela de configuração grava de forma independente.}
\ucrow{Fluxo básico}{\begin{enumerate}[nosep,leftmargin=*,label=\arabic*.]
\item A pessoa abre Configurações.
\item A pessoa escolhe a área (metas, linhas, classificações, tipos de defeito, custos, produção, áreas responsáveis).
\item O sistema mostra os valores atuais.
\item A pessoa altera e confirma.
\item O sistema valida, grava e atualiza os indicadores.
\end{enumerate}}
\ucrow{Fluxo alternativo}{\begin{enumerate}[nosep,leftmargin=*,label=A\arabic*.]
\item Valor inválido: o sistema indica o erro e não grava.
\item Falha de gravação: o sistema avisa e os valores anteriores são mantidos.
\end{enumerate}}
\end{ucard}

\Needspace{14\baselineskip}
\subsection{Assistente (gestão) e administração}

\Needspace{9\baselineskip}
\begin{ucard}{UC20}{Gerir memória do assistente}
\ucrow{Referências}{Tela: Assistente, aba Memória. Fluxo de sequência: \autoref{fig:seq08}.}
\ucrow{Descrição geral}{O caso de uso inicia-se quando a pessoa quer saber quais análises aprovadas já alimentam o assistente e em que situação estão.}
\ucrow{Atores}{Gestor; Administrador.}
\ucrow{Pré-condições}{Sessão ativa.}
\ucrow{Garantia de sucesso (pós-condições)}{Situação da memória exibida; se pedida, reindexação iniciada.}
\ucrow{Requisitos especiais}{A memória reúne só análises aprovadas. Cada análise está na memória, aguardando envio ou retirada.}
\ucrow{Fluxo básico}{\begin{enumerate}[nosep,leftmargin=*,label=\arabic*.]
\item A pessoa abre a aba Memória.
\item O sistema mostra o resumo: análises aprovadas, na memória, aguardando envio e retiradas, além da fila e do último envio.
\item A pessoa consulta a distribuição por linha.
\item Se necessário, a pessoa pede a reindexação e confirma.
\item O sistema reenvia as análises à memória e atualiza o resumo.
\end{enumerate}}
\ucrow{Fluxo alternativo}{\begin{enumerate}[nosep,leftmargin=*,label=A\arabic*.]
\item Recurso desligado: o sistema informa que a memória está desativada.
\item Memória vazia: o sistema explica que ainda não há análises aprovadas.
\item Itens com falha na fila: o resumo os destaca e a reindexação pode reenviá-los.
\end{enumerate}}
\end{ucard}

\Needspace{9\baselineskip}
\begin{ucard}{UC21}{Acompanhar desempenho do assistente}
\ucrow{Referências}{Tela: Assistente, aba Desempenho. Fluxo de sequência: \autoref{fig:seq08}.}
\ucrow{Descrição geral}{O caso de uso inicia-se quando o gestor quer saber o que as pessoas decidiram sobre as propostas do assistente.}
\ucrow{Atores}{Gestor.}
\ucrow{Pré-condições}{Decisões registradas em pré-análises (UC10).}
\ucrow{Garantia de sucesso (pós-condições)}{Indicadores de aceite, ajuste e rejeição exibidos.}
\ucrow{Requisitos especiais}{Nenhum.}
\ucrow{Fluxo básico}{\begin{enumerate}[nosep,leftmargin=*,label=\arabic*.]
\item O gestor abre a aba Desempenho.
\item O sistema calcula e exibe os indicadores de aceite, ajuste e rejeição.
\item O gestor filtra por período.
\item O sistema atualiza a visão.
\end{enumerate}}
\ucrow{Fluxo alternativo}{\begin{enumerate}[nosep,leftmargin=*,label=A\arabic*.]
\item Sem decisões no período: o sistema mostra o estado vazio.
\end{enumerate}}
\end{ucard}

\Needspace{9\baselineskip}
\begin{ucard}{UC22}{Consultar uso da plataforma}
\ucrow{Referências}{Tela: Uso da plataforma (/uso-da-plataforma).}
\ucrow{Descrição geral}{O caso de uso inicia-se quando o administrador quer ver quanto a plataforma é usada e quanto tempo ela economiza.}
\ucrow{Atores}{Administrador.}
\ucrow{Pré-condições}{Sessão ativa com perfil de administrador.}
\ucrow{Garantia de sucesso (pós-condições)}{Acessos, usuários ativos, coletas, relatórios e tempo economizado exibidos.}
\ucrow{Requisitos especiais}{O tempo economizado é o tempo manual estimado menos a duração real. Os tempos manuais são configuráveis.}
\ucrow{Fluxo básico}{\begin{enumerate}[nosep,leftmargin=*,label=\arabic*.]
\item O administrador abre Uso da plataforma.
\item O sistema exibe os indicadores do período.
\item O administrador alterna entre 30 e 90 dias.
\item O sistema atualiza o gráfico de acessos.
\end{enumerate}}
\ucrow{Fluxo alternativo}{\begin{enumerate}[nosep,leftmargin=*,label=A\arabic*.]
\item Configurar tempos manuais: o administrador informa os minutos por tarefa e o sistema recalcula a economia.
\item Sem acessos no intervalo: o sistema informa.
\end{enumerate}}
\end{ucard}

\Needspace{9\baselineskip}
\begin{ucard}{UC23}{Atender solicitações de suporte}
\ucrow{Referências}{Tela: Suporte (/suporte).}
\ucrow{Descrição geral}{O caso de uso inicia-se quando o administrador consulta as solicitações que as pessoas enviaram ao suporte.}
\ucrow{Atores}{Administrador.}
\ucrow{Pré-condições}{Solicitações registradas (UC03).}
\ucrow{Garantia de sucesso (pós-condições)}{Solicitações lidas, com assunto, mensagem e dados de quem as enviou.}
\ucrow{Requisitos especiais}{Nenhum.}
\ucrow{Fluxo básico}{\begin{enumerate}[nosep,leftmargin=*,label=\arabic*.]
\item O administrador abre a tela de Suporte.
\item O sistema lista as solicitações, da mais recente à mais antiga.
\item O administrador busca ou pagina a lista.
\item O administrador abre uma solicitação.
\item O sistema mostra o assunto, a mensagem, o perfil e a data do envio.
\end{enumerate}}
\ucrow{Fluxo alternativo}{\begin{enumerate}[nosep,leftmargin=*,label=A\arabic*.]
\item Nenhuma solicitação: o sistema mostra o estado vazio.
\item Falha ao carregar: o sistema mostra o erro e oferece tentar novamente.
\end{enumerate}}
\end{ucard}

\Needspace{9\baselineskip}
\begin{ucard}{UC24}{Configurar integração com o Teams}
\ucrow{Referências}{Tela: Integrações, Teams (/integracoes/teams). Fluxo de sequência: \autoref{fig:seq07}.}
\ucrow{Descrição geral}{O caso de uso inicia-se quando o administrador define em quais destinos do Teams os alertas devem ser entregues.}
\ucrow{Atores}{Administrador; Microsoft Teams (ator secundário).}
\ucrow{Pré-condições}{\begin{enumerate}[nosep,leftmargin=*,label=\arabic*.]
\item Sessão ativa com perfil de administrador.
\item Destino (webhook) do Teams disponível.
\end{enumerate}}
\ucrow{Garantia de sucesso (pós-condições)}{Destinos gravados; os alertas das regras associadas passam a ser entregues no Teams.}
\ucrow{Requisitos especiais}{O worker é o único contêiner que sai para fora da rede para falar com o Teams.}
\ucrow{Fluxo básico}{\begin{enumerate}[nosep,leftmargin=*,label=\arabic*.]
\item O administrador abre Integrações, Teams.
\item O sistema lista os destinos cadastrados.
\item O administrador cria ou edita um destino.
\item O sistema valida e grava o destino.
\item O worker passa a entregar os alertas ao Teams.
\end{enumerate}}
\ucrow{Fluxo alternativo}{\begin{enumerate}[nosep,leftmargin=*,label=A\arabic*.]
\item Destino inválido: o sistema indica o erro e não grava.
\item Teams indisponível na entrega: o alerta interno é mantido e a entrega ao Teams não conclui.
\end{enumerate}}
\end{ucard}

\Needspace{9\baselineskip}
\begin{ucard}{UC25}{Gerenciar usuários}
\ucrow{Referências}{Tela: Usuários (/usuarios).}
\ucrow{Descrição geral}{O caso de uso inicia-se quando o superusuário cria, altera ou desativa contas do Hanaro.}
\ucrow{Atores}{Superusuário.}
\ucrow{Pré-condições}{Sessão ativa com permissão de superusuário.}
\ucrow{Garantia de sucesso (pós-condições)}{Conta criada, alterada ou removida; o acesso passa a refletir o novo perfil.}
\ucrow{Requisitos especiais}{Toda decisão de acesso é validada também no servidor. A gravação exige o token CSRF.}
\ucrow{Fluxo básico}{\begin{enumerate}[nosep,leftmargin=*,label=\arabic*.]
\item O superusuário abre a tela de Usuários.
\item O sistema lista as contas.
\item O superusuário cria uma conta ou escolhe uma existente.
\item O superusuário define nome, usuário, perfil e senha inicial.
\item O sistema valida e grava.
\item O sistema registra a alteração na auditoria.
\end{enumerate}}
\ucrow{Fluxo alternativo}{\begin{enumerate}[nosep,leftmargin=*,label=A\arabic*.]
\item Usuário já existente: o sistema recusa e informa o conflito.
\item Remover conta: o sistema pede confirmação; a sessão da conta apagada é rejeitada nas próximas requisições.
\end{enumerate}}
\end{ucard}

\Needspace{9\baselineskip}
\begin{ucard}{UC26}{Consultar registros de atividade}
\ucrow{Referências}{Tela: Logs (/logs). Fluxo de sequência: \autoref{fig:seq13}.}
\ucrow{Descrição geral}{O caso de uso inicia-se quando o superusuário precisa auditar o que foi feito no sistema.}
\ucrow{Atores}{Superusuário.}
\ucrow{Pré-condições}{Sessão ativa com permissão de superusuário.}
\ucrow{Garantia de sucesso (pós-condições)}{Registros de atividade exibidos conforme os filtros.}
\ucrow{Requisitos especiais}{O registro é gravado pelo sistema em segundo plano; uma falha no registro não derruba a requisição da pessoa.}
\ucrow{Fluxo básico}{\begin{enumerate}[nosep,leftmargin=*,label=\arabic*.]
\item O superusuário abre a tela de Logs.
\item O sistema lista os registros, dos mais recentes aos mais antigos.
\item O superusuário filtra por intervalo, autor ou módulo, ou busca um texto.
\item O sistema atualiza a lista.
\item Se desejar, o superusuário exporta o resultado.
\end{enumerate}}
\ucrow{Fluxo alternativo}{\begin{enumerate}[nosep,leftmargin=*,label=A\arabic*.]
\item Autor inválido no filtro: o sistema indica o erro e mantém a lista anterior.
\item Registro desativado ou indisponível: o sistema informa e oferece tentar novamente.
\end{enumerate}}
\end{ucard}



\Needspace{10\baselineskip}
\section{Tela de Configurações}\label{sec:configuracoes}

A tela de Configurações reúne os dados e as regras que o Hanaro usa para classificar o que chega do GERP, comparar o scrap com as metas e montar os relatórios. Quase nada ali é calculado sozinho: o que está cadastrado na tela muda o que aparece no Dashboard, na Base de Scrap e nos relatórios. Por isso vale ler esta seção antes de alterar qualquer valor.

\subsection{Como a tela funciona}

A tela abre como um catálogo de cartões, um por seção. Clicar em um cartão abre a seção em uma janela sobre a tela, com o título e, no rodapé, os botões de salvar. Algumas regras valem para todas as seções:

\begin{itemize}[leftmargin=1.6em,itemsep=0.2em,topsep=0.3em]
\item \textbf{Quem acessa.} A tela é de analistas e gestores. Os administradores de plataforma e o superusuário são levados direto para as Execuções, porque a área deles é a operação do sistema.
\item \textbf{Alterações não salvas.} O contador \emph{“N alterações não salvas”} mostra o que ainda não foi gravado. Fechar a janela com alterações pendentes pede confirmação: \emph{Descartar} perde o que foi digitado, \emph{Continuar editando} volta para a janela.
\item \textbf{Edição ao mesmo tempo.} Se outra pessoa gravar os mesmos dados enquanto você edita, o Hanaro recusa a gravação e avisa que os dados mudaram. Use \emph{Recarregar} para ver a versão atual e refaça a alteração.
\item \textbf{Mensagens de erro.} Sem rede, sessão expirada, falta de permissão e dados recusados têm mensagens próprias. Se o aviso trouxer um \emph{código de suporte}, informe-o a quem cuida do servidor.
\end{itemize}

\begin{table}[H]
\caption{As seções da tela de Configurações.}\label{tab:config-secoes}
\guiazebra
\small
\noindent\begin{tabularx}{\textwidth}{@{}L{3.1cm}X L{3.3cm}@{}}
\toprule
\guiacabtabela
\thead{Seção} & \thead{O que controla} & \thead{Efeito} \\
\midrule
\textbf{Classificações} & Regras que traduzem os códigos do GERP em nomes e filtros do Dashboard e dos relatórios. & Próximas importações; o histórico muda ao reaplicar. \\
\textbf{Tipos de scrap} & As categorias que o analista escolhe ao classificar uma ocorrência. & Imediato, nas análises. \\
\textbf{IF Cost} & Meta mensal do ano e os custos de reparo, retrabalho e parada. & Imediato, no Dashboard. \\
\textbf{Produção} & Valor e quantidade produzidos por mês, no total ou por produto. & Imediato, no scrap em relação à produção. \\
\textbf{Programa do ano} & A ficha anual do programa de redução de scrap. & Abre os relatórios criados no ano. \\
\textbf{Cotação do dólar} & A cotação diária (PTAX), coletada sozinha do Banco Central. & Só consulta; recalcula uploads manuais. \\
\textbf{Áreas responsáveis} & A lista de áreas sugeridas ao preencher tarefas. & Sugestões nas tarefas. \\
\end{tabularx}
\end{table}

% ---------------------------------------------------------------------------
\Needspace{12\baselineskip}
\subsection{Classificações}

É a seção de maior efeito. O GERP entrega códigos (part number, organização, conta, descrição do item) e cada regra diz o que um código significa no Hanaro. Há cinco tipos de regra:

\begin{table}[H]
\caption{Tipos de regra de classificação.}\label{tab:config-regras}
\guiazebra
\small
\noindent\begin{tabularx}{\textwidth}{@{}L{3.1cm}X@{}}
\toprule
\guiacabtabela
\thead{Tipo} & \thead{O que a regra define} \\
\midrule
\textbf{Nome do item} & Um nome legível no lugar do part number na Base de Scrap e nos relatórios. Vale na hora, sem reaplicar. Campos: \emph{Part number} e \emph{Nome exibido}. \\
\textbf{Produto e divisão} & Para cada organização do GERP, os valores de \emph{Filter 1} (produto), \emph{Filter 2} e \emph{Divisão} dos filtros do Dashboard, o \emph{Código de produção} e o \emph{Mapa do painel}. Com o mapa em \emph{Sim}, a análise dos casos da organização registra a região do defeito no painel da TV. \\
\textbf{Departamento da linha} & Liga o código da linha ao departamento responsável e define se ela entra no \emph{Filtro da linha} do Dashboard. Departamentos marcados como produção aparecem sozinhos como linhas no painel Linhas, sem cadastro. \\
\textbf{Conta do IF Cost} & Diz se os lançamentos de uma conta somam no IF Cost e são comparados com as metas (\emph{Entra no IF Cost}). A regra vale para a combinação \emph{descrição da conta} + \emph{alias}. \\
\textbf{Tipo de item} & Agrupa os itens pela descrição, com texto exato ou expressão regular (maiúsculas e minúsculas são iguais). Quando várias regras combinam, vale a de menor \emph{Prioridade}. \\
\end{tabularx}
\end{table}

\subsubsection*{Como usar}

\begin{enumerate}[leftmargin=1.6em,itemsep=0.2em,topsep=0.3em]
\item Abra \emph{Classificações}. O quadro \emph{Códigos sem regra} lista os códigos que já vieram nas importações e ainda não têm regra completa, com a coluna \emph{O que falta} (nome do item ou definição de IF Cost) e a quantidade de registros afetados.
\item Filtre por \emph{Situação} (\emph{Sem regra} ou \emph{Regra incompleta}), por \emph{Falta} ou por quantidade de registros, ou busque pelo código. Comece pelos códigos com mais registros.
\item Clique em \emph{Criar regra} (ou \emph{Completar regra}) na linha do código: o formulário abre com o código preenchido. Para uma regra que não vem da lista, use \emph{Nova regra}.
\item Escolha o \emph{Tipo de regra}, preencha os campos marcados com \texttt{*} e deixe \emph{Regra ativa} ligada. Em \emph{Tipo de item} com expressão regular, teste a regra no campo \emph{Testar com uma descrição} antes de salvar.
\item Salve com \emph{Adicionar regra} (ou \emph{Salvar regra}). A regra vale para as próximas importações.
\item Para atualizar o que já foi importado, clique em \emph{Reaplicar ao histórico}. O Hanaro reclassifica todos os registros com as regras ativas e informa quantos mudaram.
\end{enumerate}

\subsubsection*{Exemplo de configuração}

Os valores abaixo são ilustrativos; os códigos reais vêm do GERP da sua fábrica.

\begin{table}[H]
\caption{Exemplos de regras (valores ilustrativos).}\label{tab:config-regras-ex}
\guiazebra
\small
\noindent\begin{tabularx}{\textwidth}{@{}L{2.7cm}X@{}}
\toprule
\guiacabtabela
\thead{Tipo} & \thead{Configuração} \\
\midrule
Nome do item & \emph{Part number} \tech{ABC123456}, \emph{Nome exibido} \tech{Placa principal 55 pol.}. A Base de Scrap passa a mostrar o nome em vez do código. \\
Conta do IF Cost & \emph{Descrição} \tech{Consumo de materiais}, \emph{Alias} \tech{E-CONSM}, \emph{Entra no IF Cost}: Sim. Os lançamentos dessa conta somam no IF Cost e entram na comparação com a meta. \\
Tipo de item & \emph{Como comparar}: Expressão regular, \emph{Expressão} \tech{\textasciicircum(?:pcba\textbar{}pcb assembly)}, \emph{Tipo de item} \tech{PCBA}, \emph{Prioridade} 10. Encontra descrições que começam com “pcba” ou “pcb assembly”. \\
Departamento da linha & \emph{Código da linha} \tech{L03}, \emph{Departamento} Montagem, \emph{Entra no Filtro da linha}: Sim. A linha passa a aparecer no filtro do Dashboard. \\
\end{tabularx}
\end{table}

\begin{remark}[atencao]
\textbf{O que não fazer em Classificações.}
\begin{itemize}[leftmargin=1.4em,itemsep=0.15em,topsep=0.3em]
\item \textbf{Não esquecer de reaplicar.} Criar, editar ou excluir uma regra não muda o que já foi importado. Sem \emph{Reaplicar ao histórico}, o Dashboard e os relatórios continuam com a classificação anterior.
\item \textbf{Não reaplicar em sequência.} A reaplicação reclassifica todos os registros. Se der erro, espere alguns minutos antes de tentar de novo, em vez de clicar várias vezes.
\item \textbf{Não usar expressão regular ampla.} Uma expressão como \tech{pcb}, sem \tech{\textasciicircum}, combina com qualquer descrição que contenha essas letras. Ancore no início, teste com descrições reais e use a prioridade para resolver sobreposições.
\item \textbf{Não duplicar regra.} Só pode existir uma regra de cada tipo por código. Se o Hanaro avisar que já existe, edite a regra existente.
\item \textbf{Não excluir para desativar.} Para parar de usar uma regra sem perdê-la, desligue \emph{Regra ativa}: regras inativas ficam guardadas, mas não classificam nada.
\item \textbf{Não desligar \emph{Entra no IF Cost} sem avisar.} A conta sai do IF Cost e da comparação com a meta, no que for importado ou reaplicado depois.
\end{itemize}
\end{remark}

% ---------------------------------------------------------------------------
\Needspace{10\baselineskip}
\subsection{Tipos de scrap}

São as categorias que o analista escolhe ao classificar uma ocorrência na análise de scrap (por exemplo, Oxidação ou Avaria de transporte). A lista é a mesma para toda a equipe.

\subsubsection*{Como usar}

\begin{enumerate}[leftmargin=1.6em,itemsep=0.2em,topsep=0.3em]
\item Em \emph{Novo tipo de scrap}, digite o \emph{Nome}. O \emph{Código} é gerado a partir do nome (letras maiúsculas, números, \texttt{\_} ou \texttt{-}) e não pode ser alterado depois.
\item Clique em \emph{Adicionar tipo}. O tipo aparece em \emph{Tipos cadastrados} e já pode ser escolhido nas análises.
\item Para mudar o nome, use \emph{Editar}. Para escondê-lo sem apagar, desligue \emph{Disponível nas análises}.
\end{enumerate}

\textbf{Exemplo.} Nome \tech{Oxidação}, código \tech{OXIDACAO}. Se a equipe deixar de usar a categoria \emph{Avaria de transporte}, desligue \emph{Disponível nas análises}: ela some das análises novas, e as análises já feitas continuam como estavam.

\begin{remark}[atencao]
\textbf{O que não fazer em Tipos de scrap.}
\begin{itemize}[leftmargin=1.4em,itemsep=0.15em,topsep=0.3em]
\item \textbf{Não excluir para esconder.} Excluir o tipo o remove das análises de vez. Para só escondê-lo, desligue \emph{Disponível nas análises}.
\item \textbf{Não escolher o código sem pensar.} O código não muda depois de criado; se estiver errado, só o nome pode ser corrigido.
\item \textbf{Não criar tipos parecidos.} “Oxidação” e “Oxidacao” viram códigos diferentes e dividem os casos em dois grupos. Procure na lista antes de adicionar.
\end{itemize}
\end{remark}

% ---------------------------------------------------------------------------
\Needspace{12\baselineskip}
\subsection{IF Cost}

Define a meta mensal de IF Cost e registra os custos de reparo, retrabalho e parada. O scrap entra sozinho, a partir das importações, nas contas marcadas como \emph{Entra no IF Cost}.

\subsubsection*{Meta do ano}

A tabela \emph{Meta por mês} guarda uma meta por mês. Acima dela, \emph{Preencher os 12 meses} gera os valores sem gravar nada, até que você clique em \emph{Salvar metas de} \emph{ano}:

\begin{itemize}[leftmargin=1.6em,itemsep=0.2em,topsep=0.3em]
\item \textbf{Reduzir de janeiro a dezembro.} Informe a \emph{Meta de janeiro} e a \emph{Meta de dezembro}; os meses do meio variam gradualmente de um valor ao outro.
\item \textbf{Mesmo valor todo mês.} Informe a \emph{Meta total do ano}; ela é dividida igualmente pelos 12 meses.
\end{itemize}

Os quadros \emph{Total planejado}, \emph{Média mensal} e \emph{Em relação a} (ano anterior) mostram o efeito da meta antes de salvar. Depois de preencher, revise a tabela: qualquer célula pode ser ajustada à mão.

\subsubsection*{Outros custos}

Em \emph{Outros custos do IF Cost}, digite reparo, retrabalho e parada de cada mês, \textbf{em dólar (US\$)}. Deixe em branco o que não tiver valor. O \emph{Total do mês} é somado na hora. Informe a \emph{Origem dos valores} (digitado manualmente, importação de planilha, ERP ou financeiro) e, se quiser, uma referência e uma observação. O \emph{Histórico de revisões} registra cada alteração.

\textbf{Exemplo.} Meta de janeiro US\$ 120~mil e de dezembro US\$ 102~mil (uma redução de 15\%): \emph{Reduzir de janeiro a dezembro}, \emph{Preencher tabela}, revisar e \emph{Salvar metas de 2026}. Em março, o financeiro envia o relatório de custos: lance o retrabalho do mês, origem \emph{Financeiro}, referência \tech{relatório financeiro de março}.

\begin{remark}[atencao]
\textbf{O que não fazer em IF Cost.}
\begin{itemize}[leftmargin=1.4em,itemsep=0.15em,topsep=0.3em]
\item \textbf{Não achar que preencher é salvar.} \emph{Preencher tabela} só monta os valores. Sem \emph{Salvar metas}, nada é gravado.
\item \textbf{Não excluir as metas do ano à toa.} \emph{Excluir metas} apaga os 12 valores e o Dashboard fica sem linha de meta nesse ano.
\item \textbf{Não lançar custos em reais.} Os campos de reparo, retrabalho e parada são em US\$.
\item \textbf{Não trocar de escopo com alterações pendentes.} Salve ou descarte antes de mudar o escopo.
\item \textbf{Não tentar lançar o scrap aqui.} Ele entra pelas importações; esta seção só recebe a meta e os outros custos.
\end{itemize}
\end{remark}

% ---------------------------------------------------------------------------
\Needspace{12\baselineskip}
\subsection{Produção}

Guarda o valor e a quantidade produzidos em cada mês. O Dashboard usa esses dados para calcular o scrap em relação à produção.

\subsubsection*{Como usar}

\begin{enumerate}[leftmargin=1.6em,itemsep=0.2em,topsep=0.3em]
\item Escolha o \emph{Escopo}: \emph{Global (todos os produtos)} ou um produto específico, e o \emph{Ano}.
\item Para cada mês, digite o \emph{Valor da produção (USD)} e a \emph{Quantidade produzida}, e uma observação se precisar. \emph{Enter} avança: valor, quantidade, próximo mês.
\item Acompanhe os meses pelas pastilhas (\emph{preenchido}, \emph{sem dados} ou \emph{valor inválido}) e pelos totais do ano.
\item Clique em \emph{Salvar dados de} \emph{ano}. Em escopo global, o quadro \emph{Detalhado por produto} mostra quanto do mês veio de cada produto.
\end{enumerate}

\textbf{Exemplo.} Escopo global, 2026, março: valor US\$ 4,2~milhões e quantidade de 38~mil unidades. Para acompanhar um produto isolado, troque o escopo, lance os meses dele e salve.

\begin{remark}[atencao]
\textbf{O que não fazer em Produção.}
\begin{itemize}[leftmargin=1.4em,itemsep=0.15em,topsep=0.3em]
\item \textbf{Não usar valores negativos ou texto.} Só números maiores ou iguais a zero; os campos inválidos ficam destacados e a gravação é recusada.
\item \textbf{Não tentar salvar com tudo vazio.} É preciso preencher pelo menos um mês.
\item \textbf{Não misturar escopos.} Lançar o total da fábrica no escopo de um produto distorce a relação entre scrap e produção.
\item \textbf{Não limpar o ano sem necessidade.} \emph{Limpar} exclui os dados e o Dashboard deixa de mostrar o scrap em relação à produção nesses meses.
\end{itemize}
\end{remark}

% ---------------------------------------------------------------------------
\Needspace{12\baselineskip}
\subsection{Programa do ano}

É a ficha anual do programa de redução de scrap. Ela abre os relatórios criados no ano, e campos vazios não aparecem no relatório. Use as setas de ano para mudar a ficha; cada ano tem a sua.

\begin{table}[H]
\caption{Blocos da ficha do programa.}\label{tab:config-ficha}
\guiazebra
\small
\noindent\begin{tabularx}{\textwidth}{@{}L{3.1cm}X@{}}
\toprule
\guiacabtabela
\thead{Bloco} & \thead{Campos} \\
\midrule
\textbf{Identificação} & Tema, período, equipe e local. \\
\textbf{KPI} & Indicador, nível atual, meta e nota da meta. Em branco, o relatório usa o valor da base e das metas, atualizado na emissão; preencha só para substituir um valor. \emph{Valor manual} e \emph{Voltar ao automático} alternam entre os dois modos. \\
\textbf{Contexto e resultados} & Problema de negócio, ideias de ruptura, resultados qualitativos e pontos críticos; um item por linha. \\
\textbf{Plano de execução} & Etapas com período e o que é feito, até o limite indicado na tela. \\
\textbf{Organização} & Patrocinador, responsável, líder e as equipes de tempo integral e parcial; um nome por linha. \\
\end{tabularx}
\end{table}

\textbf{Exemplo.} Tema \tech{Material Scrap TASK – MS IF Cost redução de 15\%}, período \tech{Jan–Dez}, KPI \tech{IF Cost de scrap de material (US\$)}, com nível atual e meta deixados em branco para virem da base. Salve com \emph{Salvar ficha}; o rodapé passa a indicar a versão salva.

\begin{remark}[atencao]
\textbf{O que não fazer na ficha.}
\begin{itemize}[leftmargin=1.4em,itemsep=0.15em,topsep=0.3em]
\item \textbf{Não digitar o KPI por hábito.} Um valor manual congela o número e deixa de acompanhar a base. Use \emph{Voltar ao automático} quando a razão para o valor manual acabar.
\item \textbf{Não editar a ficha de outro ano por engano.} Confira o ano no topo antes de salvar.
\item \textbf{Não ignorar o aviso de conflito.} Se a ficha mudou em outra sessão, recarregue antes de salvar, ou a alteração é recusada.
\end{itemize}
\end{remark}

% ---------------------------------------------------------------------------
\Needspace{10\baselineskip}
\subsection{Cotação do dólar}

Mostra a cotação de venda do dólar (PTAX), coletada sozinha do Banco Central. O sistema confere novidades a cada hora. Os uploads usam a cotação do dia em que o arquivo chegou; em fins de semana e feriados, a do dia útil anterior. Nada nesta seção é digitado: ela serve para conferir.

\begin{itemize}[leftmargin=1.6em,itemsep=0.2em,topsep=0.3em]
\item O aviso de estado diz se a coleta está normal, atrasada (traz a data da última cotação recebida) ou ainda vazia.
\item O gráfico e a tabela do ano mostram a cotação de cada dia e quantos dias úteis ficaram sem cotação (feriado ou falha na coleta).
\item \emph{Recalcular dados carregados} refaz o valor em US\$ dos \textbf{uploads manuais} já feitos, com as cotações de agora. Dados vindos da automação não mudam.
\end{itemize}

\begin{remark}[atencao]
\textbf{O que não fazer na cotação.}
\begin{itemize}[leftmargin=1.4em,itemsep=0.15em,topsep=0.3em]
\item \textbf{Não esperar corrigir a cotação à mão.} Ela não é editável; se a coleta estiver atrasada por muito tempo, avise quem cuida do servidor.
\item \textbf{Não recalcular sem motivo.} Use \emph{Recalcular dados carregados} depois que uma cotação faltante chegar; o recálculo altera valores de uploads já conferidos.
\end{itemize}
\end{remark}

% ---------------------------------------------------------------------------
\Needspace{8\baselineskip}
\subsection{Áreas responsáveis}

Cadastra as áreas sugeridas ao preencher o campo de área de uma tarefa (por exemplo, Qualidade). Todos veem as sugestões; só o superusuário cadastra e remove áreas.

\begin{itemize}[leftmargin=1.6em,itemsep=0.2em,topsep=0.3em]
\item Digite o \emph{Nome da área} e clique em \emph{Cadastrar área}. Uma área repetida é recusada.
\item \emph{Remover} tira a área das sugestões; as tarefas existentes mantêm o valor que já tinham.
\end{itemize}

\begin{remark}[dica]
Padronize o nome das áreas (“Qualidade”, não “qualidade” nem “Qualid.”): as sugestões servem justamente para que as tarefas fiquem agrupáveis nos relatórios.
\end{remark}


% Gerado por gen_agent.py. Não editar à mão.
\Needspace{10\baselineskip}
\section{Agente de IA e RAG da consulta}\label{sec:agente-rag}

O servidor de IA (\texttt{agent-service}) faz dois tipos de trabalho. O primeiro roda em segundo plano, como \textbf{job}: a pré-análise de uma ocorrência (revisão da análise e propostas de tarefas) e as propostas de tarefas de um relatório. O segundo é a \textbf{consulta}, que responde na hora a uma pergunta em texto livre. Em todos, o modelo só adapta texto: tudo o que não precisa de modelo é código, e toda saída passa por uma crítica em código antes de chegar à pessoa. A IA propõe, a pessoa decide.

Nos diagramas, as caixas rosa chamam um modelo (de linguagem ou de embeddings), as caixas areia são código puro, os losangos são decisões e as setas tracejadas são leituras ou escritas no banco.

\Needspace{10\baselineskip}
\subsection{Revisão de análises: o grafo do LangGraph}

A pré-análise é um grafo do LangGraph. O estado carrega só valores JSON e o grafo grava um checkpoint no PostgreSQL a cada passo, com o identificador do job como \texttt{thread}. Se o worker cair, o job é retomado do passo que foi interrompido, sem repetir os anteriores.

\begin{figure}[H]
\centering
\includegraphics[width=\textwidth,height=\dimexpr\textheight-3.2cm\relax,keepaspectratio]{img/diagramas/agente-g1}
\caption{Máquina de estados do grafo de revisão de análises (LangGraph).}\label{fig:agente-g1}
\end{figure}

\begin{table}[H]
\caption{Nós do grafo de revisão.}\label{tab:agente-nos}
\guiazebra
\small
\noindent\begin{tabularx}{\textwidth}{@{}L{3.6cm}X L{1.9cm}@{}}
\toprule
\guiacabtabela
\thead{Nó} & \thead{O que faz} & \thead{Modelo} \\
\midrule
\texttt{facts} & Monta os fatos da ocorrência (item, linha, período, valor) e o gráfico base de comparação. & Não \\
\texttt{retrieve} & Gera o embedding do caso, busca no pgvector os casos parecidos e a reincidência e junta exemplos que analistas já aceitaram. & Embedding \\
\texttt{grade} & Decide a rota: a (há relatório concluído do mesmo item), b (há casos parecidos) ou c (nenhum). & Não \\
\texttt{compose\_without\_precedent} & Rota c: título e descrição só com os fatos, sem causa, e o checklist dos 4M. & Não \\
\texttt{adapt\_review} & Propõe os campos da revisão (título, descrição, causa, tipo de defeito), cada um com fontes. & Forte \\
\texttt{adapt\_tasks} & Propõe tarefas a partir das tarefas dos casos parecidos; se nenhum caso tem tarefas, não chama o modelo. & Forte \\
\texttt{investigate} & Opcional (desligado por padrão): até 5 consultas de leitura sobre os casos recuperados; o resultado só vai ao registro técnico. & Forte \\
\texttt{critique} & Recusa o que não cita fonte válida, traz número ou data fora das entradas, usa valor inválido ou estoura limite, e qualquer causa na rota c. A certeza da causa é fixada em código como hipótese. & Não \\
\texttt{finalize} & Monta e valida o resultado final. & Não \\
\end{tabularx}
\end{table}

As três rotas mudam o que o grafo pode afirmar. Na \textbf{rota a} já existe um relatório concluído do mesmo item (reincidência); na \textbf{rota b} há casos parecidos; na \textbf{rota c} não há nenhum, e então o texto só repete os fatos, sem causa. Se a crítica recusa algo, o grafo volta uma única vez a \texttt{adapt\_review} com a lista de erros; na segunda passagem, o que continua recusado é descartado e aparece como descartado no resultado.

\Needspace{10\baselineskip}
\subsection{Criação de tarefas para um relatório}

Quando a pessoa pede propostas de tarefas para um relatório que já existe (\texttt{POST /v1/task-proposals}), o serviço segue uma sequência curta, sem LangGraph e sem checkpoint: se o job falhar, a nova tentativa recomeça do início. A regra é a mesma da pré-análise: só se propõe uma tarefa se ela se apoia em tarefas de casos parecidos. Sem caso parecido, ou com casos que não têm tarefas, o serviço não chama o modelo e devolve uma pergunta ao analista. Tarefas que o relatório já tem são removidas, e cada proposta recebe uma chave estável (\texttt{proposal\_key}), calculada do título e das fontes.

\begin{figure}[H]
\centering
\includegraphics[width=\textwidth,height=\dimexpr\textheight-3.2cm\relax,keepaspectratio]{img/diagramas/agente-g2}
\caption{Máquina de estados da criação de tarefas para um relatório.}\label{fig:agente-g2}
\end{figure}

\Needspace{10\baselineskip}
\subsection{Ciclo de vida do job}

Os jobs da pré-análise e das propostas de tarefas usam a mesma fila no PostgreSQL. A API aceita o pedido com uma chave de idempotência e responde \texttt{202}. O worker pega o job mais antigo com \texttt{FOR UPDATE SKIP LOCKED}, então vários workers não disputam o mesmo. Falhas que podem se repetir devolvem o job à fila com espera crescente de 30 segundos, 2 minutos e 10 minutos, até 3 tentativas. Um job cujo worker parou de responder é retomado por outro; se parou na última tentativa, é marcado como falho.

\begin{figure}[H]
\centering
\includegraphics[width=\textwidth,height=\dimexpr\textheight-3.2cm\relax,keepaspectratio]{img/diagramas/agente-g3}
\caption{Ciclo de vida de um job do agente.}\label{fig:agente-g3}
\end{figure}

\Needspace{10\baselineskip}
\subsection{RAG da consulta}

A consulta é um RAG: antes de escrever, o sistema recupera casos da \textbf{memória}, que reúne as análises aprovadas. A resposta é montada só com o que foi recuperado. O fluxo é fixo, e o modelo não escolhe o que fazer: reescrever a pergunta, buscar, decidir se há caso bom o bastante, escrever e verificar. A consulta responde dentro da própria requisição (\texttt{POST /v1/chat}) e tem poucas vagas simultâneas; se estiverem ocupadas, o serviço responde que está ocupado.

\Needspace{8\baselineskip}
\subsubsection*{Indexação e pgvector}

Cada análise aprovada vira uma linha em \texttt{cases}, com os campos usados em filtros e o snapshot completo em JSONB. O texto do caso (título, defeito, família 4M, descrição, causas, porquês e tarefas) é calculado em código e vira um vetor, guardado em \texttt{case\_embeddings} numa coluna \texttt{vector} do pgvector. Um índice HNSW com distância de cosseno (\texttt{m = 16}, \texttt{ef\_construction = 64}) acelera a busca dos vizinhos mais próximos. O hash do texto evita gerar de novo o vetor de um caso que não mudou, e a busca só considera vetores do modelo de embeddings em uso.

\begin{figure}[H]
\centering
\includegraphics[width=\textwidth,height=\dimexpr\textheight-3.2cm\relax,keepaspectratio]{img/diagramas/agente-ri}
\caption{Indexação: da análise aprovada ao vetor no pgvector.}\label{fig:agente-ri}
\end{figure}

\Needspace{8\baselineskip}
\subsubsection*{Busca híbrida, fusão e reranking}

A pergunta é reescrita por um modelo rápido, que devolve o texto de busca e filtros (item, linha, defeito, período). Um filtro só vale se existir algum caso com aquele valor exato; os que o modelo inventa são descartados. Em seguida, duas buscas rodam com os mesmos filtros e trazem 30 casos cada:

\begin{itemize}
\item \textbf{Léxica (BM25):} pontua as palavras, e mantém códigos de item, linhas e códigos de relatório como termos exatos.
\item \textbf{Semântica (pgvector):} compara o vetor da pergunta com os vetores dos casos pela distância de cosseno.
\end{itemize}

As duas listas são fundidas por \textbf{RRF} (Reciprocal Rank Fusion, com k = 60): cada caso soma 1/(60 + posição) em cada lista, o que dispensa pesos e normalização. Se o servidor de embeddings estiver fora, a consulta segue só pela busca léxica.

O \textbf{reranking} é um segundo estágio opcional, desligado por padrão: um cross-encoder, servido por uma rota \texttt{/rerank} do servidor de modelos, pontua cada par de pergunta e caso e reordena os candidatos. Só se liga depois que a avaliação de recuperação (\texttt{eval/retrieval.py}) mostrar ganho, porque um reranker também sobe casos parecidos que enganam o redator. Se ele falhar, vale a ordem da fusão.

Um caso só chega ao redator se estiver perto o bastante. \textbf{Com reranker}, a nota absoluta precisa ser de pelo menos 0,3. \textbf{Sem reranker}, basta a similaridade de cosseno ser de pelo menos 0,5 ou as palavras da pergunta cobrirem ao menos metade do caso. No máximo 6 casos seguem; se nenhum passa, a resposta é \texttt{no\_precedent} e o redator nem é chamado.

\begin{figure}[H]
\centering
\includegraphics[width=\textwidth,height=\dimexpr\textheight-3.2cm\relax,keepaspectratio]{img/diagramas/agente-rs}
\caption{Busca híbrida da consulta, com pgvector, BM25, fusão e reranking.}\label{fig:agente-rs}
\end{figure}

\Needspace{8\baselineskip}
\subsubsection*{Resposta e grounding}

O redator (modelo forte) recebe cada caso resumido ao que responde a pergunta e devolve um resumo e afirmações, cada uma com os casos que a sustentam. O \textbf{grounding} é feito em código, depois do modelo:

\begin{itemize}
\item a afirmação cita ao menos um caso que foi de fato oferecido ao redator (citar um caso que não foi oferecido invalida a fonte);
\item nenhum número ou data aparece sem estar nos casos ou na conversa, e o texto cabe em 1000 caracteres;
\item o resumo também é conferido contra números inventados.
\end{itemize}

Se alguma afirmação é recusada, o redator tenta uma segunda vez, recebendo a lista de erros. Ao final, as afirmações recusadas ficam em \texttt{discarded}, com o motivo, e o resultado tem um de três estados: \texttt{answered} (ao menos uma afirmação válida, com as fontes citadas), \texttt{no\_valid\_answer} (há casos, mas nenhuma afirmação se sustenta; só as fontes são mostradas) e \texttt{no\_precedent}.

\begin{figure}[H]
\centering
\includegraphics[width=\textwidth,height=\dimexpr\textheight-3.2cm\relax,keepaspectratio]{img/diagramas/agente-rb}
\caption{Resposta da consulta e verificação de grounding.}\label{fig:agente-rb}
\end{figure}

\Needspace{8\baselineskip}
\subsubsection*{A busca da pré-análise é outra}

A pré-análise não usa BM25 nem reranking. Ela procura casos que compartilham campos com a ocorrência (mesmo item, tipo, linha, defeito e família de causa) e os mais próximos pelo pgvector, e pontua tudo com pesos: item 5, tipo 2, linha 2, defeito 2, família de causa 1, texto até 2 e boa prática 1. O resultado perde metade do valor a cada 12 meses e é multiplicado pela utilidade do caso, que aprende com as decisões das pessoas; um caso marcado como "não é parecido" perde metade da nota. Um caso só vira precedente se tiver relevância mínima de 3,0 e algo em comum além do contexto (mesmo item, mesmo defeito ou texto parecido).



\section{Componentes}

Cinco diagramas, do mais amplo ao mais detalhado: o sistema inteiro e, depois, o
interior de cada serviço. As páginas estão giradas para aproveitar a folha.

\begin{figure}[H]
\centering
\includegraphics[width=\dimexpr\textheight-3.2cm\relax,height=\textwidth,keepaspectratio,angle=90]{img/diagramas/componentes-sistema}
\caption{Componentes do sistema.}\label{fig:componentes-sistema}
\end{figure}

\begin{figure}[H]
\centering
\includegraphics[width=\dimexpr\textheight-3.2cm\relax,height=\textwidth,keepaspectratio,angle=90]{img/diagramas/componentes-frontend}
\caption{Componentes do front-end.}\label{fig:componentes-frontend}
\end{figure}

\begin{figure}[H]
\centering
\includegraphics[width=\dimexpr\textheight-3.2cm\relax,height=\textwidth,keepaspectratio,angle=90]{img/diagramas/componentes-backend}
\caption{Componentes do back-end (API).}\label{fig:componentes-backend}
\end{figure}

\begin{figure}[H]
\centering
\includegraphics[width=\dimexpr\textheight-3.2cm\relax,height=\textwidth,keepaspectratio,angle=90]{img/diagramas/componentes-worker}
\caption{Componentes do worker.}\label{fig:componentes-worker}
\end{figure}

\begin{figure}[H]
\centering
\includegraphics[width=\dimexpr\textheight-3.2cm\relax,height=\textwidth,keepaspectratio,angle=90]{img/diagramas/componentes-agent-service}
\caption{Componentes do servidor de IA.}\label{fig:componentes-agent-service}
\end{figure}

\Needspace{10\baselineskip}
\section{Implantação}\label{sec:implantacao-uml}

O diagrama de implantação da UML mostra \emph{onde cada parte do Hanaro roda}. A \autoref{fig:implantacao-uml} usa três elementos. Cada caixa em relevo é um \textbf{nó}: um dispositivo (\texttt{device}), como um servidor, ou um ambiente de execução (\texttt{executionEnvironment}), como um contêiner ou um navegador. O retângulo com ícone dentro de um nó é um \textbf{artefato}, o que foi implantado nele. As linhas são os \textbf{caminhos de comunicação}, com a seta indo de quem chama para quem responde. A figura da Parte~I mostra o mesmo cenário de forma simplificada; aqui estão todos os contêineres, as redes e os sentidos das chamadas.

\begin{table}[H]
\caption{Nós e artefatos do diagrama de implantação.}\label{tab:impl-nos}
\guiazebra
\small
\noindent\begin{tabularx}{\textwidth}{@{}L{3.3cm}X@{}}
\toprule
\guiacabtabela
\thead{Nó} & \thead{Artefatos e observações} \\
\midrule
\textbf{Navegador} & A aplicação web (Angular), servida pelo nginx. \\
\textbf{Automação GERP} & O processo \texttt{hanaro-automation}, em Python: lê o arquivo TSV exportado do GERP e envia a carga ao Hanaro com uma chave de API. \\
\textbf{Servidor do Hanaro} & Os contêineres do Hanaro, em Docker Compose. Só o nginx publica porta; a rede \texttt{data} é interna, sem saída para fora. \\
\textbf{nginx} & O artefato \texttt{hanaro-web}. Rede \texttt{proxy}; porta 4203; sistema de arquivos somente leitura. \\
\textbf{backend} & O artefato \texttt{hanaro-backend} (a API). Redes \texttt{proxy} e \texttt{data}; a porta 8000 só existe dentro da rede; aplica as migrações do banco ao iniciar. \\
\textbf{worker} & A mesma imagem do backend, com o comando do Taskiq. Redes \texttt{data} e \texttt{proxy}; é o único contêiner que também sai para fora. \\
\textbf{postgres e redis} & O banco de dados e o Redis (fila, sessões e cache do Dashboard), na rede \texttt{data}, cada um com seu volume. \\
\textbf{Servidor de IA} & Opcional e, em geral, em outra máquina. Reúne o compose do agente (\texttt{agent-api} na porta 8200, \texttt{agent-worker} e \texttt{agent-db}, um PostgreSQL com pgvector), o compose da tradução (\texttt{translation-service}, porta 8100) e o servidor de modelos (Ollama ou equivalente, com API compatível com a da OpenAI). \\
\textbf{Banco Central e Teams} & Sistemas externos: o Banco Central fornece a cotação PTAX (Olinda) e o Teams recebe os cartões de alerta. Só o worker fala com eles. \\
\end{tabularx}
\end{table}

\begin{figure}[H]
\centering
\includegraphics[width=\dimexpr\textheight-3.2cm\relax,height=\textwidth,keepaspectratio,angle=90]{img/diagramas/implantacao}
\caption{Diagrama de implantação do Hanaro.}\label{fig:implantacao-uml}
\end{figure}

\begin{table}[H]
\caption{Caminhos de comunicação do diagrama de implantação.}\label{tab:impl-caminhos}
\guiazebra
\small
\noindent\begin{tabularx}{\textwidth}{@{}L{3.3cm}L{4.2cm}X@{}}
\toprule
\guiacabtabela
\thead{Caminho} & \thead{De, para} & \thead{Observação} \\
\midrule
HTTP, porta 4203 & Navegador, nginx & A tela e a API usam o mesmo endereço. \\
HTTP com chave de API & Automação, nginx & A API não publica porta própria, então a carga entra pelo nginx. \\
\texttt{/api} e \texttt{/health} & nginx, backend & O nginx serve a tela e encaminha a API. \\
SQL, fila e sessões & backend e worker, postgres e redis & Os dois contêineres usam o banco e o Redis. \\
HTTP, chat & backend, \texttt{agent-api} & A consulta do analista. \\
HTTP, jobs & worker, \texttt{agent-api} & O worker entrega os pedidos de pré-análise e busca os resultados. O agente nunca chama o Hanaro. \\
HTTP, tradução & worker, \texttt{translation-service} & Textos já mascarados. \\
HTTP & agente e tradução, servidor de modelos & Os dois serviços de IA chamam o servidor de modelos. \\
HTTPS & worker, Banco Central e Teams & Cotação PTAX e cartões de alerta. \\
\end{tabularx}
\end{table}

\begin{remark}[info]
\textbf{Para TI.} As portas 8200 do agente e 8100 da tradução devem ficar fechadas, exceto para o servidor do Hanaro. O serviço \texttt{populate-demo}, que só existe para preencher uma instalação nova com dados de demonstração, não aparece no desenho.
\end{remark}

% ---------------------------------------------------------------------------
\Needspace{10\baselineskip}
\section{Pacotes}\label{sec:pacotes}

O diagrama de pacotes mostra como o código está organizado e quem pode depender de quem. Cada caixa com aba é um \textbf{pacote}, uma pasta de código; pacotes pequenos dentro de outro estão contidos nele. A seta vai de quem \emph{usa} para quem é \emph{usado}, e a seta de duas pontas indica uso nos dois sentidos. As regras desenhadas são as que o repositório confere em uma checagem automática.

\subsection{Backend}

O código da API fica em seis camadas e usa uma infraestrutura compartilhada (\autoref{fig:pacotes-backend}).

A checagem automática impede quatro coisas: \texttt{controller} importar \texttt{repository}; \texttt{services} importar \texttt{controller}; \texttt{repository} importar \texttt{controller}, \texttt{services} ou \texttt{support}; e \texttt{models} ou \texttt{utils} importarem qualquer camada acima deles. Dois pontos merecem atenção ao mudar o código: \texttt{services} e \texttt{support} se usam nos dois sentidos, e \texttt{repository} existe só para alguns módulos, porque os serviços de scrap e de governança consultam o banco diretamente pela sessão da infraestrutura.

\begin{table}[H]
\caption{Pacotes do backend.}\label{tab:pacotes-backend}
\guiazebra
\small
\noindent\begin{tabularx}{\textwidth}{@{}L{2.8cm}X L{3.4cm}@{}}
\toprule
\guiacabtabela
\thead{Pacote} & \thead{O que contém} & \thead{Pode usar} \\
\midrule
\texttt{controller} & Rotas e validação da entrada, por domínio (acesso, scrap, governança, uso da plataforma, suporte e administração). & \texttt{services}, \texttt{support}, \texttt{models} \\
\texttt{services} & Casos de uso e regras de negócio, por domínio. & \texttt{repository}, \texttt{support}, \texttt{models}, \texttt{utils} \\
\texttt{support} & Dependências, exceções, tarefas do worker e renderizadores, por domínio. & \texttt{models} e, em alguns módulos, \texttt{services} \\
\texttt{repository} & Acesso a dados de alguns módulos (chaves de API, scrap, limite de taxa, perfis de limite e usuários). & \texttt{models}, \texttt{utils} \\
\texttt{models} & Tabelas e esquemas de dados, por domínio. & \texttt{utils} \\
\texttt{utils} & Funções puras. & Nada da aplicação \\
\texttt{infrastructure} & Sessão do banco, cache, fila Taskiq, autenticação, limite de taxa, segurança, registro (logging), middleware e fábrica da aplicação. & Todas as camadas a usam \\
\end{tabularx}
\end{table}

\begin{figure}[H]
\centering
\includegraphics[width=\dimexpr\textheight-3.2cm\relax,height=\textwidth,keepaspectratio,angle=90]{img/diagramas/pacotes-backend}
\caption{Pacotes do backend.}\label{fig:pacotes-backend}
\end{figure}

\subsection{Aplicação web}

A aplicação Angular tem um pacote por funcionalidade navegável dentro de \texttt{modules} (\autoref{fig:pacotes-frontend}).

As regras são estas: \texttt{modules} usa \texttt{core} e \texttt{shared}; \texttt{core} pode usar \texttt{shared}; \texttt{shared} não importa de \texttt{core} nem de \texttt{modules}. Uma funcionalidade só usa outra pelos arquivos \texttt{*.public-api.ts} dela, e uma checagem automática confere essas fronteiras. Regras e modelos exclusivos de uma tela ficam na própria funcionalidade; o código só sobe para \texttt{core} ou \texttt{shared} quando há consumidores em mais de uma delas.

\begin{table}[H]
\caption{Pacotes da aplicação web.}\label{tab:pacotes-frontend}
\guiazebra
\small
\noindent\begin{tabularx}{\textwidth}{@{}L{2.8cm}X@{}}
\toprule
\guiacabtabela
\thead{Pacote} & \thead{O que contém} \\
\midrule
\texttt{app} & Rotas, configuração da aplicação e os layouts. \\
\texttt{modules} & Uma pasta por funcionalidade: dashboard, executions, scrap-base, reports, governance, assistant, settings, auth, users, profile, admin-logs, platform-usage, support e tv. \\
\texttt{core} & Responsabilidades globais, com uma única instância: sessão e autenticação, configuração, cliente HTTP, idioma, áreas responsáveis, estrutura da tela e tema. \\
\texttt{shared} & Componentes, gráficos, diretivas, pipes, utilitários, estilos e textos usados por mais de uma funcionalidade. \\
\end{tabularx}
\end{table}

\begin{figure}[H]
\centering
\includegraphics[width=\dimexpr\textheight-3.2cm\relax,height=\textwidth,keepaspectratio,angle=90]{img/diagramas/pacotes-frontend}
\caption{Pacotes da aplicação web.}\label{fig:pacotes-frontend}
\end{figure}


\section{Fluxos de sequência}

Treze fluxos, na ordem em que o sistema é usado: acesso, entrada de dados, análise, relatório e auditoria. Cada diagrama mostra quem chama quem e em que ordem.

\subsection{Login e sessão}

A sessão fica no Redis e o navegador só guarda o identificador, em cookie \texttt{HttpOnly}. O token CSRF volta no corpo da resposta e deve ser enviado nas escritas.

\begin{figure}[H]
\centering
\includegraphics[width=\textwidth,height=\dimexpr\textheight-4.5cm\relax,keepaspectratio]{img/diagramas/seq01}
\caption{Login e sessão.}\label{fig:seq01}
\end{figure}

\subsection{Ingestão automática do GERP}

A automação (hoje com arquivo local, no futuro com o robô) converte o arquivo do GERP no JSON canônico. A API só valida a forma e enfileira; quem reconcilia e grava é o worker. Repetir o mesmo arquivo na mesma janela é idempotente.

\begin{figure}[H]
\centering
\includegraphics[width=\textwidth,height=\dimexpr\textheight-4.5cm\relax,keepaspectratio]{img/diagramas/seq02}
\caption{Ingestão automática do GERP.}\label{fig:seq02}
\end{figure}

\subsection{Envio manual e acompanhamento da execução}

O mesmo caminho do diagrama 2, mas com o arquivo vindo do navegador. A cotação usada é a do dia de chegada informado (hoje em Manaus por padrão); sem cotação nesse dia, vale a última disponível.

\begin{figure}[H]
\centering
\includegraphics[width=\textwidth,height=\dimexpr\textheight-4.5cm\relax,keepaspectratio]{img/diagramas/seq03}
\caption{Envio manual e acompanhamento da execução.}\label{fig:seq03}
\end{figure}

\subsection{Cotação do dólar (PTAX)}

É um laço do worker, sem ação do usuário. A tela de Configurações apenas mostra que a coleta está funcionando e onde há lacunas. Não existe cotação manual.

\begin{figure}[H]
\centering
\includegraphics[width=\dimexpr\textheight-4.5cm\relax,height=\textwidth,keepaspectratio,angle=90]{img/diagramas/seq04}
\caption{Cotação do dólar (PTAX).}\label{fig:seq04}
\end{figure}

\subsection{Leitura do dashboard com cache}

O dashboard lê uma projeção já agregada, não as transações. A chave do cache inclui a revisão da projeção; uma ingestão ou a mudança de meta troca a revisão. Se o Redis falhar, a resposta sai do banco.

\begin{figure}[H]
\centering
\includegraphics[width=\dimexpr\textheight-4.5cm\relax,height=\textwidth,keepaspectratio,angle=90]{img/diagramas/seq05}
\caption{Leitura do dashboard com cache.}\label{fig:seq05}
\end{figure}

\subsection{Análise de uma ocorrência}

O analista abre uma ocorrência da Base de scrap, preenche a análise como rascunho e a finaliza. Há uma análise ativa por ocorrência, e as gravações usam a versão esperada para não pisar em outra edição.

\begin{figure}[H]
\centering
\includegraphics[width=\textwidth,height=\dimexpr\textheight-4.5cm\relax,keepaspectratio]{img/diagramas/seq06}
\caption{Análise de uma ocorrência.}\label{fig:seq06}
\end{figure}

\subsection{Regra de alerta, alerta interno e Teams}

O worker avalia as regras a cada ciclo (30 s). Quando um nível dispara, o fato vira um evento na caixa de saída, gravado na mesma transação da avaliação. Depois o evento é despachado: cria o alerta interno e enfileira o cartão do Teams, que outro laço (15 s) entrega com nova tentativa em caso de falha.

\begin{figure}[H]
\centering
\includegraphics[width=\dimexpr\textheight-4.5cm\relax,height=\textwidth,keepaspectratio,angle=90]{img/diagramas/seq07}
\caption{Regra de alerta, alerta interno e Teams.}\label{fig:seq07}
\end{figure}

\subsection{Pré-análise do assistente e aceite}

O Hanaro nunca espera o agente dentro de uma requisição ou de uma transação de alerta: ele grava o pedido e o worker o entrega depois. O agente nunca chama o Hanaro; o worker busca o resultado. O aceite é transacional e idempotente.

\begin{figure}[H]
\centering
\includegraphics[width=\textwidth,height=\dimexpr\textheight-4.5cm\relax,keepaspectratio]{img/diagramas/seq08}
\caption{Pré-análise do assistente e aceite.}\label{fig:seq08}
\end{figure}

\subsection{Chat de consulta do assistente}

O chat responde só sobre os casos da memória e cita as fontes. A conversa é do usuário que a abriu.

\begin{figure}[H]
\centering
\includegraphics[width=\dimexpr\textheight-4.5cm\relax,height=\textwidth,keepaspectratio,angle=90]{img/diagramas/seq09}
\caption{Chat de consulta do assistente.}\label{fig:seq09}
\end{figure}

\subsection{Relatório: da edição à publicação}

Publicar sela uma fotografia dos dados (snapshot com hash) e cria uma versão imutável numerada. A publicação aceita uma chave de idempotência para que um clique repetido não gere duas versões.

\begin{figure}[H]
\centering
\includegraphics[width=\textwidth,height=\dimexpr\textheight-4.5cm\relax,keepaspectratio]{img/diagramas/seq10}
\caption{Relatório: da edição à publicação.}\label{fig:seq10}
\end{figure}

\subsection{Exportação do relatório}

A exportação é assíncrona: o pedido vira um job, o evento sai pela caixa de saída e o worker renderiza o arquivo. O job tem uma reserva (lease) para que só um worker o processe e para recuperar o que ficou parado.

\begin{figure}[H]
\centering
\includegraphics[width=\textwidth,height=\dimexpr\textheight-4.5cm\relax,keepaspectratio]{img/diagramas/seq11}
\caption{Exportação do relatório.}\label{fig:seq11}
\end{figure}

\subsection{Tradução de um plano de ação}

O Hanaro protege números e códigos, mascara o texto e só então chama o serviço de tradução. Só um job fala com o serviço por vez; os outros ficam na fila e a tela mostra a posição.

\begin{figure}[H]
\centering
\includegraphics[width=\textwidth,height=\dimexpr\textheight-4.5cm\relax,keepaspectratio]{img/diagramas/seq12}
\caption{Tradução de um plano de ação.}\label{fig:seq12}
\end{figure}

\subsection{Auditoria de atividade}

Um middleware registra quem fez o quê nas escritas das áreas de acesso, scrap, alertas e relatórios. O registro vem depois da resposta e uma falha ao gravar não derruba a requisição. A exportação é só para superusuário.

\begin{figure}[H]
\centering
\includegraphics[width=\dimexpr\textheight-4.5cm\relax,height=\textwidth,keepaspectratio,angle=90]{img/diagramas/seq13}
\caption{Auditoria de atividade.}\label{fig:seq13}
\end{figure}

\section{Modelo conceitual de dados}

O modelo conceitual está dividido em cinco recortes para manter a leitura. Algumas entidades, como Usuário, Caso e Plano de ação, aparecem em mais de um recorte, sempre com o mesmo nome. Uma relação não aparece desenhada: o \emph{Usuário} se relaciona consigo mesmo, pois um usuário (gestor) pode gerenciar outros usuários.

\begin{figure}[H]
\centering
\includegraphics[width=\dimexpr\textheight-4.5cm\relax,height=\textwidth,keepaspectratio,angle=90]{img/diagramas/merA}
\caption{Modelo conceitual, recorte 1.}\label{fig:merA}
\end{figure}

\begin{figure}[H]
\centering
\includegraphics[width=\dimexpr\textheight-4.5cm\relax,height=\textwidth,keepaspectratio,angle=90]{img/diagramas/merB}
\caption{Modelo conceitual, recorte 2.}\label{fig:merB}
\end{figure}

\begin{figure}[H]
\centering
\includegraphics[width=\textwidth,height=\dimexpr\textheight-4.5cm\relax,keepaspectratio]{img/diagramas/merC1}
\caption{Modelo conceitual, recorte 3: relatório publicado.}\label{fig:merC1}
\end{figure}

\begin{figure}[H]
\centering
\includegraphics[width=\textwidth,height=\dimexpr\textheight-4.5cm\relax,keepaspectratio]{img/diagramas/merC2}
\caption{Modelo conceitual, recorte 4: tarefas e casos.}\label{fig:merC2}
\end{figure}

\begin{figure}[H]
\centering
\includegraphics[width=\textwidth,height=\dimexpr\textheight-4.5cm\relax,keepaspectratio]{img/diagramas/merD}
\caption{Modelo conceitual, recorte 5.}\label{fig:merD}
\end{figure}

\fimdecapitulo
